% id: 50-11
% book: 베이직쎈 미적분1 · 03 미분계수와 도함수
% section: 유형 022 평균변화율과 미분계수의 기하적 의미
% figure: tikz:fig-50-11
% answer: 
% answer_source: 
% variant_level: 0 (원본 전사)
% 이 파일은 build.mjs 가 items.json 에서 만든다. 고칠 때는 items.json 을 고친다.
\smash{\raisebox{4.5mm}{\dmkichul{베이직쎈 미적분1 50쪽 11번}}}
\noindent\begin{minipage}[t]{0.58\linewidth}\vspace{0pt}\raggedright 오른쪽 그림과 같이 $y$축에 대하여 대칭인 이차함수 $y=f(x)$의 그래프 위의 점~$\pt{A}(2{,}\mkern3mu\mathopen{}f(2))$에 대하여 직선~$\pt{OA}$의 기울기가 $\dfrac{1}{2}$이다. 함수 $f(x)$에서 $x$의 값이 $-2$에서 $0$까지 변할 때의 평균변화율은? \dmcondfill{$\pt{O}$는 원점이다.}{}\end{minipage}\hfill\begin{minipage}[t]{0.4\linewidth}\vspace{0pt}\centering \resizebox{\ifdim\width>\linewidth\linewidth\else\width\fi}{!}{% 50-11(50쪽) — y축에 대하여 대칭인 이차함수 y=f(x) 의 그래프와 그래프 위의 점 A(2, f(2)) 를 지나는 직선 OA. 스캔 화질이 낮아 TikZ 로 다시 그림.
% 원문에 있는 것만: 포물선(꼭짓점이 원점) · 원점을 지나는 직선 · A 에서 x축으로 내린 점선 · 라벨 y=f(x), A, O, 2, x, y. 눈금·f(2) 값은 원문에 없어 적지 않는다.
\begin{tikzpicture}[x=0.36cm, y=0.40cm, line width=0.5pt]
  \draw[->, >=stealth, line width=0.7pt] (-4.4,0) -- (4.7,0) node[below=1pt] {$x$};
  \draw[->, >=stealth, line width=0.7pt] (0,-1.4) -- (0,4.5) node[left=1pt] {$y$};
  \node[below left=-2pt and -3pt, font=\small] at (0,0) {$\pt{O}$};
  % y축에 대하여 대칭인 이차함수 (꼭짓점 = 원점)
  \draw[line width=0.6pt, domain=-3.72:3.72, samples=90, smooth] plot (\x, {0.25*\x*\x});
  % 직선 OA (A 의 좌우로 연장)
  \draw[line width=0.5pt] (-2.4,-1.2) -- (4.4,2.2);
  % A 에서 x축에 내린 점선
  \draw[densely dashed, line width=0.4pt] (2,1) -- (2,0);
  \node[above left=-1pt and -3pt, font=\small] at (2,1) {$\pt{A}$};
  \node[below=1pt, font=\small] at (2,0) {$2$};
  \node[anchor=west, font=\small] at (2.55,3.95) {$y=f(x)$};
\end{tikzpicture}
}\end{minipage}\par
\dmchoicesauto{\rule[-3ex]{0pt}{8ex}}{$-1$}{$-\dfrac{5}{6}$}{$-\dfrac{2}{3}$}{$-\dfrac{1}{2}$}{$-\dfrac{1}{3}$}
